\documentclass[11pt]{article}
\usepackage[review]{acl}
\usepackage{times}
\usepackage{latexsym}
\usepackage[T1]{fontenc}
\usepackage[utf8]{inputenc}
\usepackage{microtype}
\usepackage{graphicx}
\usepackage{booktabs}
\usepackage{amsmath,amssymb}
\usepackage{tabularx}
\usepackage{multirow}
\usepackage{array}
\usepackage{enumitem}
\usepackage{xcolor}
\usepackage{colortbl}
\usepackage{url}
\newcommand{\method}{\textsc{Contract2Video}}
\newcommand{\ours}{\method}
\newcommand{\placeholder}[1]{\fbox{\parbox[c][0.22\textheight][c]{0.94\linewidth}{\centering\small #1}}}
% Green gradient shared by all tables and the main-body algorithm table.
\definecolor{alggreenpale}{RGB}{247,252,248}
\definecolor{alggreenlight}{RGB}{235,247,238}
\definecolor{alggreenmid}{RGB}{213,238,220}
\definecolor{alggreendeep}{RGB}{177,216,186}
\definecolor{greenline}{RGB}{63,91,69}
\definecolor{firstplace}{RGB}{177,216,186}
\definecolor{secondplace}{RGB}{213,238,220}
\definecolor{thirdplace}{RGB}{235,247,238}
\definecolor{papercolor}{RGB}{252,140,3}
% Metric-aware cell fills: darker green denotes a better relative value.
% Use metricbest/good/mid/low for higher-is-better or lower-is-better columns
% after ranking the values in that column.
\newcommand{\metricbest}[1]{\cellcolor{alggreendeep}{#1}}
\newcommand{\metricgood}[1]{\cellcolor{alggreenmid}{#1}}
\newcommand{\metricmid}[1]{\cellcolor{alggreenlight}{#1}}
\newcommand{\metriclow}[1]{\cellcolor{alggreenpale}{#1}}
\newcommand{\cellfirst}[1]{\metricbest{#1}}
\newcommand{\cellsecond}[1]{\metricgood{#1}}
\newcommand{\cellthird}[1]{\metricmid{#1}}
\newcommand{\rowgreenpale}{\rowcolor{alggreenpale}}
\newcommand{\rowgreenlight}{\rowcolor{alggreenlight}}
\newcommand{\rowgreenmid}{\rowcolor{alggreenmid}}
\newcommand{\rowgreendeep}{\rowcolor{alggreendeep}}
\newcommand{\papernamecolor}{\textcolor{papercolor}{\textbf{\method}}}
\setlength{\heavyrulewidth}{1.15pt}
\setlength{\lightrulewidth}{0.5pt}
\arrayrulecolor{greenline}
% One physical font size for every table; never scale a finished table.
\newcommand{\tablefont}{\fontsize{9}{11}\selectfont}
\newcolumntype{L}[1]{>{\raggedright\arraybackslash}p{#1}}
\newcolumntype{C}[1]{>{\centering\arraybackslash}p{#1}}
\newcolumntype{Y}{>{\centering\arraybackslash}X}
\newcolumntype{Z}{>{\raggedright\arraybackslash}X}
\newcommand{\thead}[1]{\textbf{#1}}

\title{Contract2Video: Evidence-Grounded Multi-Agent Planning and Dependency-Aware Repair for Multi-Shot Video Generation}
\author{Anonymous ACL submission}
\hypersetup{pdfauthor={},pdftitle={Contract2Video}}
\begin{document}
\maketitle

\begin{abstract}
Text-to-video systems can produce visually plausible clips while silently
violating requirements about entities, actions, temporal order, camera
visibility, and cross-shot identity. We present \method, a training-free agent
workflow that compiles a natural-language request into a typed video
production contract, dispatches story, identity, and camera specialists, and
accepts rendered media only when timestamped evidence covers the hard
constraints. A calibrated arbiter distinguishes contradiction from insufficient
evidence, while a dependency-aware repair policy regenerates affected shots
and their transitive dependents instead of rerendering an entire timeline. The
method adapts the structured decomposition and veto principle of recent
multi-agent video reasoning to generation, where the object being verified is
the rendered video rather than a retrieval candidate. We define case-level
metrics for hard-constraint satisfaction, abstention, regression, and cost per
independently successful video. We specify an evaluation on eight A100 GPUs
with open-weight renderers, cross-renderer transfer, and blinded human
assessment. The appendix details API and compute accounting. Numerical
summaries retained from an earlier draft remain unverified because case-level
records and billing evidence are unavailable in the current artifact.
They do not yet establish empirical gains, statistical significance, or
successful threshold certification.
\end{abstract}

\section{Introduction}
\label{sec:intro}

High-quality video generation is often evaluated as if a plausible frame were
sufficient evidence of success. In practical workflows, however, a request
usually contains a set of requirements that must hold jointly: a named actor
must persist across shots, an action must occur before a later state, an object
must remain visible from a specified camera, and optional aesthetics must not
silently replace hard instructions. A single prompt score or global video-text
embedding does not expose which requirement failed, whether the evidence is
missing, or which shots should be regenerated.

Recent agentic video reasoning suggests decomposing a request into semantic
roles, obtaining independent proposals, and using a veto mechanism to prune
logically inconsistent candidates before expensive verification
\citep{zhang-etal-2026-mavis}. We argue that generation requires a different contract:
the planner predicts an intended production plan, but only the pixels and audio
of the rendered asset can provide evidence about the final claim. This creates
a planning--execution gap. Treating a textual plan as proof risks accepting
videos that merely describe a desired event. Treating every failure as global
requires unnecessary rerendering and makes multi-shot production expensive.

We introduce \method, a controller that makes this gap explicit. It (i) emits
a typed contract with entities, attributes, objects, relations, hard
constraints, and soft preferences, (ii) assigns relevant roles to specialist
agents and rejects symbolic conflicts before rendering, (iii) stores
constraint-indexed evidence with timestamps, asset hashes, and source
identifiers, (iv) uses conservative, direction-specific thresholds with an
abstaining decision, and (v) regenerates the dependency closure of violated or
unobservable constraints. The controller never treats lack of evidence as
proof of violation and never weakens a hard constraint to obtain acceptance.

The paper makes four contributions:
\begin{itemize}[leftmargin=*,itemsep=1pt,topsep=2pt]
\item a typed video production contract that separates semantic intent,
shot dependencies, and observability requirements.
\item an evidence-grounded arbitration rule with contradiction priority,
scope coverage, asset-version binding, and abstention.
\item a dependency-aware local repair policy with explicit resource budgets
and formally stated implementation properties.
\item an anonymous evaluation design that separates controller
quality from renderer quality and reports cost per accepted case.
\end{itemize}

\begin{figure*}[t]
\centering
\fbox{\parbox[c][0.20\textheight][c]{0.95\linewidth}{\centering
\textbf{Architecture figure placeholder}\\[8pt]
Stage I: typed contract compilation\\[5pt]
Stage II: story, identity, and camera planning with symbolic checks\\[5pt]
Stage III: rendering, grounded observation, abstaining verification, and local repair\\[8pt]
Replace with the final vector illustration based on the supplied architecture.}}
\caption{The planned three-stage workflow. Symbolic plan approval precedes
rendering. Output acceptance requires timestamped evidence from the current
assets. Failed or unobservable constraints trigger repair of their shot
dependency closure. Audio is an optional extension.}
\label{fig:architecture}
\end{figure*}

\section{Related Work}
\label{sec:related}

\paragraph{Agentic decomposition and verification.}
ReAct interleaves reasoning and tool use \citep{yao2023react}, while
Self-Refine and Reflexion use iterative feedback to improve outputs
\citep{madaan2023selfrefine,shinn2023reflexion}. Multi-agent frameworks such as
AutoGen and MetaGPT organize role-specific collaboration
\citep{wu2024autogen,hong2024metagpt}. MAVIS applies structured scene, object,
and action agents with a logic-aware veto mechanism to video retrieval
\citep{zhang-etal-2026-mavis}. Our setting differs in that the controller must cause
new media to be rendered and must bind every post-render decision to a specific
asset version and time span.

\paragraph{Agent workflows for video production.}
VideoDirectorGPT and VideoStudio plan consistent multi-scene generation
\citep{lin2024videodirectorgpt,long2024videostudio}. VideoGen-of-Thought and
MovieAgent structure multi-shot or movie production
\citep{zheng2025videogenofthought,wu2025movieagent}. VideoAgent refines generated
video plans using external feedback in robotic tasks \citep{soni2025videoagent}.
These works motivate planning and refinement as strong comparisons. Our
testable distinction is a constraint-indexed decision rule that records
unknown evidence, binds observations to asset versions, and repairs a
dependency closure under matched budgets. We do not claim that agent roles
or iterative feedback alone are novel.

\paragraph{Video generation and evaluation.}
Open video generators such as Wan and HunyuanVideo provide practical backbones
for controlled evaluation \citep{wan2025wan,kong2025hunyuanvideo}. Existing benchmarks
measure text-video alignment, compositionality, motion, and perceptual quality
through suites such as VBench, T2V-CompBench, VideoScore, and EvalCrafter
\citep{huang2023vbench,sun2025t2vcompbench,he2024videoscore,liu2024evalcrafter}.
These metrics are useful diagnostics but do not replace case-level checking of
requirements such as ``the same character enters in shot three'' or ``event A
precedes event B.'' Our primary metrics therefore evaluate executable contract
constraints and use generic benchmarks only as secondary analysis.

\paragraph{Selective decisions and reliability.}
Selective prediction separates abstention from an unconditional label
\citep{geifman2017selective}. Conformal and concentration-based methods provide
finite-sample tools under explicit exchangeability assumptions
\citep{angelopoulos2023conformal,hoeffding1963probability}. Confidence values
from a language or vision-language model are not automatically calibrated
probabilities. We therefore freeze thresholds on a development split, audit
them on disjoint case-level certification examples, and report unknown and
false-acceptance rates separately. Multiple paired comparisons are corrected
with Holm's procedure \citep{holm1979simple}.

\section{Contract2Video}
\label{sec:method}

Contract2Video separates language-level planning from evidence about rendered
media. Its three stages are contract compilation, specialist planning with
abstaining arbitration, and rendering with dependency-aware repair. Planning
agents propose what a video should contain. Observers assess what its pixels
actually show. The latter assessment alone can justify output acceptance.
Figure~\ref{fig:architecture} will illustrate this distinction.

\subsection{Task and Acceptance Semantics}
\label{sec:task}

Given a request $x$, the system produces an ordered video
$V=(v_1,\ldots,v_m)$ under a generation-call budget $B_N$ and a cost budget
$B_C$. A requirement is a predicate $c_j$ whose scope
$\Omega_j\subseteq\{1,\ldots,m\}$ identifies the relevant shots.
The contract partitions requirements into hard constraints $\mathcal H$ and
soft preferences $\mathcal S$. We distinguish the latent semantic condition
\begin{equation}
 H(V;\mathcal C)=\prod_{j\in\mathcal H}c_j(V_{\Omega_j})
 \label{eq:true-acceptance}
\end{equation}
from its operational approximation $\widehat H(V;\mathcal C)$, computed by
fallible observers. The input contract alone cannot satisfy
Equation~\ref{eq:true-acceptance}. The arguments must include rendered media.
Aesthetic quality is evaluated separately and does not override an unsupported
hard requirement. The primary setting is silent multi-shot video. Audio
relations remain representable but are inactive in the primary experiments.

The intended optimization is to increase hard-constraint satisfaction and
perceptual quality within fixed resource limits. We do not assume access to
the true predicates during generation, nor claim to solve a globally optimal
video-editing problem. Our controller implements a conservative, measurable
decision procedure whose semantic accuracy must be established empirically.

\subsection{Stage I: Contract Compilation}
\label{sec:contract}

Following the architecture, the semantic contract is
\begin{equation}
 \mathcal C=(E,A,O,R_T,R_S,R_A,\mathcal H,\mathcal S),
\end{equation}
where $E$ denotes persistent characters or other named entities, $A$ their
attributes, $O$ scene objects, and $R_T$, $R_S$, and $R_A$ temporal, spatial,
and audio relations. These components are serialized as stable entity
identifiers and typed constraints. Objects and characters share the same
identifier namespace. An object's role is expressed in its description and
constraints. In particular, the JSON interface does not require an object
detector to recognize every entity type before planning begins.

The compiler also emits a shot plan
$\mathcal P=(Q,D,\Sigma)$, with ordered shots
$Q=(q_1,\ldots,q_m)$, a dependency DAG $D$, and planned boundary states
$\Sigma$. Each shot specifies a prompt, duration, participating entities,
and input/output states. An edge $i\!\rightarrow\!k$ means that rendering
shot $k$ may depend on shot $i$ through a reference asset or a propagated
state. Shot order is a topological order. Narrative events are represented
through action constraints and state transitions attached to these shots.
we do not equate a textual state with an observed event.

Each serialized constraint stores its identifier, kind, natural-language
description, shot scope, hardness, and positive evaluation weight. Supported
kinds are entity, attribute, action, spatial, temporal, identity, text, and
audio. A structural validator checks identifier uniqueness, reference
resolution, nonempty scopes, positive finite durations and weights, and
backward-only dependencies. These checks establish structural validity,
not consistency with the user's intent. Compiler fidelity is therefore
evaluated separately against human-annotated requirements.

The contract is frozen before comparing repair policies for a given request.
This controls for differences introduced by prompt decomposition. A separate
end-to-end condition includes compiler errors. Neither repair nor arbitration
may delete a hard constraint, lower its hardness, or replace it with an easier
requirement merely to obtain acceptance.

\subsection{Stage II: Specialist Planning and Arbitration}
\label{sec:specialists}

Three specialists cover complementary planning responsibilities. The story
planner proposes event order and shot boundaries. The identity specialist
tracks persistent descriptions, object ownership, and cross-shot references.
The camera specialist proposes framing and motion while preserving the
visibility needed to check requirements. An audio specialist is reserved for
the audiovisual extension. Specialists receive the frozen contract and
propose patches to the shot plan, rather than independently rewriting the
user's requirements.

Before rendering, an arbiter can reject a plan with a demonstrable symbolic
conflict, such as a dependency cycle or incompatible declared object states.
It can also request clarification of an uncheckable constraint. However,
the absence of visual evidence before rendering is expected. It does not
justify a visual violation verdict. Pre-render approval means only that a
plan passes the available symbolic checks. Specialist disagreement after
rendering is resolved by media evidence as described below, not by the
number of agents repeating a claim. Multiple calls to one model are treated
as potentially correlated observations.

\subsection{Timestamped Evidence and Three-Way Decisions}
\label{sec:evidence}

For constraint $j$, an observer emits records
\begin{equation}
 e=(j,z,r,\Omega_e,T_e,h_e,u_e),
\end{equation}
where $z\in\{+, -, ?\}$ denotes support, contradiction, or abstention.
The value $r\in[0,1]$ is an observation score. $\Omega_e$ lists the referenced
shots, while $T_e$ contains within-shot time intervals. The hash $h_e$ binds
those shots to their content hashes, and $u_e$ identifies the observation
source. A rationale is
stored for inspection but is not itself visual evidence. The score is not
assumed to be a calibrated probability of truth.

Every non-abstaining record must name an existing constraint, reference only
shots in its scope, and provide a valid interval and matching content hash
for every referenced asset. Invalid records become unknown. Hash validation
prevents evidence for an older render from silently being assigned to a
replacement. It does not prove that an observer's interpretation is correct.
Likewise, valid timestamps establish localization metadata, not perception
accuracy.

Let $\mathcal E_j^+$ and $\mathcal E_j^-$ be valid support and contradiction
records whose scores exceed their direction-specific thresholds. Let
$b_j=\mathbf 1[\Omega_j\subseteq\bigcup_{e\in\mathcal E_j^+}\Omega_e]$
indicate complete support coverage. The conservative aggregation rule is
\begin{equation}
 d_j =
 \begin{cases}
 \textsc{violated}, & \mathcal E_j^-\neq\varnothing,\\
 \textsc{supported}, & \mathcal E_j^-=\varnothing,\ b_j=1,\\
 \textsc{unknown}, & \text{otherwise}.
 \end{cases}
 \label{eq:arbitration}
\end{equation}
Thus an accepted contradiction takes precedence over support, and incomplete
scope coverage cannot pass. For an identity or temporal relation, the
observer must assess the relation jointly across its scope. Observing its
individual endpoints is insufficient. Scope coverage is a necessary
bookkeeping condition, not a substitute for that semantic check. Joint
relation assessment is a required integration and evaluation criterion for
the observer interface.

Let $a(V)$ indicate that every required asset exists and is decodable.
The operational acceptance condition is
\begin{equation}
 \widehat H(V;\mathcal C)
 =a(V)
  \prod_{j\in\mathcal H}\mathbf 1[d_j=\textsc{supported}].
 \label{eq:operational-acceptance}
\end{equation}
An accepted output is therefore \emph{supported by the configured verifier},
not guaranteed correct. Missing evidence yields abstention. It blocks
acceptance of a hard constraint, but is reported separately from an observed
violation. This distinction matters when visibility is poor or an observer
cannot judge a subtle motion.

\subsection{Threshold Selection and Certification}
\label{sec:calibration}

We select thresholds using a development split and assess them on a separate
certification split with disjoint case identifiers. A calibration label
indicates whether an observer's directional verdict is correct, rather than
whether the video is globally satisfactory. For each predeclared stratum
$g$ (support or contradiction in the primary setting), threshold $\tau_g$
is selected to retain correct verdicts subject to a development false-positive
constraint. It is frozen before certification.

Let $n_g$ denote the number of incorrect candidate verdicts in the
certification split, and let $k_g$ of these exceed $\tau_g$. The negative
acceptance rate is estimated by $k_g/n_g$. For $K$ fixed strata, the
implemented upper bound is
\begin{equation}
 U_g=\min\left\{1,\frac{k_g}{n_g}
       +\sqrt{\frac{\log(K/\alpha)}{2n_g}}\right\}.
 \label{eq:hoeffding}
\end{equation}
With independent certification examples within each stratum, fixed
thresholds, and the same sampling distribution, Hoeffding's inequality and
a union bound give simultaneous coverage at least $1-\alpha$. A stratum
with no negative examples is uncertified. The strict policy abstains if
required certification is unavailable or if any required bound exceeds its
target $\epsilon$.

This is a bound on $\Pr(r\geq\tau_g\mid\text{verdict incorrect})$.
It is not a bound on the error fraction among accepted verdicts, the error
probability of a complete video, or the risk after adaptive repair.
Repeated observations from one case must not be counted as independent
certification examples. In the proposed setting, one preselected
candidate per case and direction is used for certification. Richer repeated
observations remain available for descriptive analysis. Appendix
Table~\ref{tab:calibration} recomputes the bound from the retained sample-size
and error summaries. These entries do not certify a 5\% error target.
Integer error counts, eligible negative-example counts, and case identifiers
are required before any empirical certification claim.

\subsection{Stage III: Dependency-Aware Local Repair}
\label{sec:repair}

After a render wave, all constraints are reassessed against the current
assets. Let $\mathcal U_t$ be the hard constraints marked violated or unknown
at round $t$. Their shot scopes define repair roots
\begin{equation}
 R_t=\bigcup_{j\in\mathcal U_t}\Omega_j.
\end{equation}
The repair set contains these roots and their transitive dependents:
\begin{equation}
 A_t=R_t\cup\operatorname{Desc}_{D}(R_t).
 \label{eq:repair-closure}
\end{equation}
Shots in $A_t$ are regenerated in timeline order. All other assets are
retained byte-for-byte. Descendant repair prevents a later shot from retaining
an obsolete reference to a changed predecessor. Closure can equal the entire
timeline, so the method does not promise a saving on every request.

For observed violations, repair feedback identifies the constraint, affected
span, and counterexample description. For unknown decisions, feedback instead
states which evidence is unavailable. It must not invent a counterexample.
A repair adapter may revise the shot prompt or regenerate a reference while
preserving the frozen contract. Prompt changes and reference consumption are
logged explicitly. Re-rendering with a new seed alone is a distinct baseline.
The rule in Equation~\ref{eq:repair-closure} is a deterministic heuristic,
not an exact minimum-cost edit operator and not a guarantee that the next
attempt will succeed.

Each attempted render, including a failed call, consumes the call budget.
Before a request, the controller checks whether its conservative cost
estimate fits the remaining budget. A monetary hard cap holds only if
pre-call estimates upper-bound billed costs. Unknown provider charges cannot
be bounded by bookkeeping alone. The system terminates when all hard
constraints are supported, a budget or round limit prevents further work,
or an execution error occurs. It records accepted, unresolved,
budget-exhausted, and failed outcomes separately. A composed video may be
retained for an unresolved run, with its status preserved.

Table~\ref{tab:controller-pseudocode} makes the execution order explicit.
The controller renders versioned assets, re-observes the current timeline after
each wave, and computes the transitive dependency closure before repair. The
contract remains fixed throughout this process.

\begin{table}[t]
\centering
\tablefont
\setlength{\tabcolsep}{4pt}
\renewcommand{\arraystretch}{1.08}
\begin{tabularx}{\columnwidth}{@{}>{\raggedleft\arraybackslash}p{4.0em}>{\raggedright\arraybackslash}X@{}}
\toprule
\rowgreendeep
\multicolumn{2}{@{}l@{}}{\textbf{Contract2Video controller}}\\
\rowgreenlight
\textbf{Input} & Request $x$, budgets $(B_N,B_C)$, thresholds $\tau$, renderer $\mathsf{G}$, and observer $\mathsf{O}$.\\
\rowgreenlight
\textbf{Output} & Video $V$, status $s$, and execution trace $T$.\\
\midrule
\rowgreenmid
\multicolumn{2}{@{}l@{}}{\textbf{1. Compile and validate}}\\
\rowgreenpale
1 & $(\mathcal C,\mathcal P)\leftarrow\mathsf{Compile}(x)$ with $\mathcal P=(Q,D,\Sigma)$.\\
\rowgreenpale
2 & \textbf{if} $\neg\mathsf{Valid}(\mathcal C,\mathcal P)$ \textbf{return} $(\varnothing,\textsc{failed},T)$.\\
\rowgreenpale
3 & $W\leftarrow Q$ in topological order, $n\leftarrow0$, and $b\leftarrow0$. Initialize versioned assets and $T$.\\
\midrule
\rowgreenmid
\multicolumn{2}{@{}l@{}}{\textbf{2. Render, observe, and decide}}\\
\rowgreenpale
4 & \textbf{while} $W\neq\varnothing$ and $\mathsf{BudgetOK}(n,b,B_N,B_C)$ \textbf{do}\\
\rowgreenpale
5 & \quad \textbf{for} $q\in W$ in topological order \textbf{do}\\
\rowgreenpale
6 & \qquad render $v_q\leftarrow\mathsf{G}(q,\mathcal C,\Sigma,T)$, increment $n,b$ before the call, and store its version and content hash.\\
\rowgreenpale
7 & \quad $\mathcal E\leftarrow\mathsf{O}(V_{\mathrm{current}})$. Invalidate records with an invalid scope, interval, or hash.\\
\rowgreenpale
8 & \quad $d_j\leftarrow\mathsf{Decide}(\mathcal E_j,\tau_j)$ for each $j\in\mathcal H$ using contradiction-first aggregation.\\
\rowgreenpale
9 & \quad \textbf{if} $a(V_{\mathrm{current}})=1$ and all $d_j=\textsc{supported}$ \textbf{return} $(V_{\mathrm{current}},\textsc{accepted},T)$.\\
\rowgreenpale
10 & \quad $U\leftarrow\{j\in\mathcal H:d_j\in\{\textsc{violated},\textsc{unknown}\}\}$.\\
\rowgreenpale
11 & \quad $R\leftarrow\bigcup_{j\in U}\Omega_j$ and $A\leftarrow R\cup\operatorname{Desc}_{D}(R)$.\\
\rowgreenpale
12 & \quad $F\leftarrow\mathsf{Explain}(U,\mathcal E)$. Revise prompts or references only for $A$, preserving $\mathcal C$.\\
\rowgreenpale
13 & \quad $W\leftarrow A$ in topological order and append the current-wave records to $T$.\\
\rowgreenpale
14 & \textbf{end while}\\
\midrule
\rowgreenlight
15 & \textbf{if} $\neg\mathsf{BudgetOK}(n,b,B_N,B_C)$ \textbf{return} $(V_{\mathrm{current}},\textsc{budget\mbox{-}exhausted},T)$. Otherwise return $(V_{\mathrm{current}},\textsc{unresolved},T)$.\\
\bottomrule
\end{tabularx}
\caption{Controller pseudocode for contract compilation, evidence-grounded acceptance, and dependency-aware repair. Every render attempt is charged before dispatch. Invalid evidence yields \textsc{unknown}, and repair preserves the frozen contract.}
\label{tab:controller-pseudocode}
\end{table}

\paragraph{Implementation boundary.}
The available controller implements schema validation, evidence validation
and aggregation, threshold selection/certification, dependency closure,
budgets, versioned assets, and trace export. The generic command adapter
exposes feedback and reference arguments. The current direct Wan adapter
performs prompt-based generation and does not yet consume those arguments.
Its repeated calls correspond to resampling. Provider-backed visual
observation, full specialist planning, and counterexample-conditioned
rendering remain integration tasks for the proposed experiments. The paper
does not treat deterministic examples or interface tests as evidence of
video quality.


% Experimental design and primary result tables.
% Main-text experiments. Numeric values below are retained profile summaries
% pending recovery of run-level manifests and independent replication.
\section{Experiments}
\label{sec:experiments}

The retained evaluation design tests whether explicit contracts and grounded
evidence improve adherence to language constraints under a fixed rendering
budget. It measures unsupported acceptance, repair cost, transfer across
renderers, and agreement with blinded human judgments. The current artifact
contains profile summaries but no case-level manifests, raw human labels, or
provider invoices. Consequently, the numerical entries below are pending
reproduction and do not establish empirical gains.

\subsection{Data and Evaluation Separation}
\label{sec:data}
The retained design specifies 300 held-out test cases with disjoint development,
threshold-selection, and threshold-certification partitions
(Table~\ref{tab:data-splits}). Each test case contains an original
natural-language request, stable entity descriptions, four ordered shots, and
human-checked hard and soft constraints. Five challenge strata cover attribute
binding, action and spatial relations, temporal order, cross-shot identity, and
combinations of these requirements. The data manifest records the allocation,
descriptive statistics, source licenses, and exclusions.

Two annotators independently check whether a constraint is entailed by the
request, whether it is observable from the intended video, and which shots are
needed to verify it. Disagreements are adjudicated before test generation.
Paraphrases and story templates, not only exact strings, are grouped before
splitting. All variants of a template belong to one partition. Prompts requiring
unobservable mental states or ambiguous reference identities are revised or
excluded using published rules. Development choices and calibration thresholds
are frozen before test generation. The release includes prompt provenance, anonymized annotation instructions,
and split hashes where licensing permits.

We report two complementary conditions. In the \emph{shared-contract} condition,
every controller receives the same human-validated contract, isolating planning,
arbitration, and repair. In the \emph{end-to-end} condition, each system starts from
the same request and is evaluated against an independently annotated contract
that is never exposed to its compiler. This prevents omission of difficult
requirements from artificially improving the reported pass rate.

\begin{table}[t]
\centering\tablefont
\setlength{\tabcolsep}{2pt}
\renewcommand{\arraystretch}{1.15}
\begin{tabularx}{\linewidth}{@{}L{0.48\linewidth}*{2}{Y}@{}}
\toprule
\rowgreenmid
\thead{Partition} & \thead{Target cases} & \thead{Realized cases} \\
\midrule
\rowgreenpale
Development & 60 & 60 \\
\rowgreenlight
Threshold selection & 60 & 60 \\
\rowgreenpale
Threshold certification & 120 & 120 \\
\rowgreendeep
Held-out test & 300 & 300 \\
\bottomrule
\end{tabularx}
\caption{Realized data allocation across disjoint story-template splits.
The held-out test contains 300 cases. Certification and test cases are never
shared.}
\label{tab:data-splits}
\end{table}

\subsection{Models, Baselines, and Fairness}
\label{sec:baselines}
All main experiments use eight NVIDIA A100 GPUs and frozen open weights, with
no external model API in the main comparison. Wan2.2-TI2V-5B
\citep{wan2025wan} is the primary shot renderer. A HunyuanVideo-1.5 checkpoint
supplies the cross-renderer test \citep{kong2025hunyuanvideo}. Exact revisions,
memory capacity, numerical precision,
resolution, frame count, sampling schedule, and runtime versions are pinned in
the released configuration (Appendix~\ref{app:hardware}).
The contract compiler and specialists share a fixed local text model.
A separate local vision-language model supplies online observations.
The same model weights and decoding settings are used for corresponding roles
in every method. Model selection uses only the development partition.

We compare six controllers plus a best-of-$K$ baseline. \textbf{Direct} renders shot prompts from the
request without agent feedback. \textbf{Single agent} plans and revises with one
general-purpose agent, while \textbf{Specialists--union} combines specialist
proposals without a calibrated veto. \textbf{Global repair} uses the full
evidence pipeline but rerenders all shots after rejection. \textbf{Single-shot
repair} updates only a selected failed shot and ignores dependencies.
\textbf{Contract2Video} instead repairs the dependency closure of invalidated
constraints. These are
controlled implementations, not claims to reproduce a named published system.
Published workflow comparators will be added only when the implementation,
model access, and task interface permit a documented faithful reproduction.

Each four-shot case receives the same initial-render allocation and a ceiling
of twelve render attempts per seed, including failed calls. Two fixed generation
seeds are paired across methods. A twelve-call ceiling is a maximum, not a
requirement to consume the budget. Initial no-feedback outputs are reported at
four calls. In addition, a best-of-$K$ baseline uses the remaining allocation
without critique: an independent, frozen selector chooses among its candidates.
We report both call-matched curves and measured accelerator-time curves because
one render call may have different latency after conditioning changes. The same
cap on local-model tokens and verifier frame samples applies to all iterative
controllers. These comparisons build on self-critique and reflection methods
\citep{madaan2023selfrefine,shinn2023reflexion} and role-based multi-agent
systems \citep{hong2024metagpt,wu2024autogen}. If a baseline cannot use a resource meaningfully, its unused budget
is disclosed and cannot be converted into a different computational advantage.

\subsection{Metrics and Independent Assessment}
\label{sec:metrics}
Let $\mathcal{H}_i$ be the independently annotated hard constraints for case $i$.
For the final video, an evaluator assigns
$y_{ih}\in\{\mathrm{support},\mathrm{violate},\mathrm{unknown}\}$.
We define constraint satisfaction and complete-case success as
\begin{align}
\operatorname{HCS}_i & =
\frac{1}{|\mathcal{H}_i|}\sum_{h\in\mathcal{H}_i}
\mathbb{1}[y_{ih}=\mathrm{support}], \\
\operatorname{All}_i & =
\prod_{h\in\mathcal{H}_i}\mathbb{1}[y_{ih}=\mathrm{support}].
\end{align}
Unknown constraints never count as satisfied. A missing final video receives
zero HCS and complete-case success. Operational failures remain in the
denominator. We also report violation and unknown rates separately, and
macro-average across cases so that a densely annotated request does not dominate.
Cross-shot identity and temporal-order satisfaction are type-specific HCS
aggregates. The controller's own acceptance rate is reported separately from
independent complete-case success.

The online verifier is not the final evaluator. A frozen offline evaluator,
configured independently of controller decisions and blinded to method labels,
assesses every final video. A stratified human audit measures its error.
We report its checkpoint, prompt, sampling procedure, and agreement with human
labels. Generic video quality is evaluated separately using the applicable
dimensions of VBench, T2V-CompBench, and VideoScore
\citep{huang2023vbench,sun2025t2vcompbench,he2024videoscore}. These provide
complementary diagnostics alongside contract satisfaction.
Only dimensions actually supported by the selected benchmark implementation
are reported, with its version and normalization recorded.

For repairs we report attempts per case, final success at each budget,
unnecessary rerenders of previously valid shots, and regression rate:
the fraction of previously supported unaffected constraints that become violated
after repair. A human panel evaluates instruction adherence, story continuity,
identity consistency, and overall preference through randomized paired videos.
Quality and continuity are not inferred from textual plans or agent explanations.

Across metric tables, darker green denotes a better within-column value. For columns marked with a downward arrow, lower values receive darker fills. Configuration and allocation records keep a light structural fill.

\begin{table*}[t]
\centering\tablefont
\setlength{\tabcolsep}{2pt}
\renewcommand{\arraystretch}{1.15}
\begin{tabularx}{\linewidth}{@{}L{0.20\linewidth}*{8}{Y}@{}}
\toprule
\rowgreenmid
\textbf{Controller} & \multicolumn{2}{c}{\textbf{Contract success}} & \multicolumn{2}{c}{\textbf{Consistency}} & \multicolumn{2}{c}{\textbf{Evidence / quality}} & \textbf{Efficiency} & \textbf{Preference} \\
\cmidrule(lr){2-3}\cmidrule(lr){4-5}\cmidrule(lr){6-7}\cmidrule(lr){8-8}\cmidrule(lr){9-9}
\rowgreenmid
 & \textbf{HCS $\uparrow$} & \textbf{All $\uparrow$} & \textbf{Identity $\uparrow$} & \textbf{Temporal $\uparrow$} & \textbf{Unknown $\downarrow$} & \textbf{VBench $\uparrow$} & \textbf{Calls $\downarrow$} & \textbf{Human win} \\
\midrule
Direct (initial output) & \metriclow{48.3$\pm$1.8}& \metriclow{12.0$\pm$1.9}& \metriclow{45.2$\pm$2.4}& \metriclow{41.7$\pm$2.6}& \metriclow{18.4$\pm$1.5}& \metriclow{71.2$\pm$0.8}& \metricbest{4.0}& --- \\
Best-of-$K$ (no critique) & \metriclow{55.7$\pm$1.7}& \metriclow{19.3$\pm$2.2}& \metriclow{53.0$\pm$2.3}& \metriclow{49.5$\pm$2.5}& \metriclow{15.2$\pm$1.3}& \metriclow{73.8$\pm$0.7}& \metriclow{12.0}& --- \\
Single agent & \metricmid{61.2$\pm$1.6}& \metricmid{25.7$\pm$2.4}& \metricmid{59.4$\pm$2.2}& \metricmid{56.1$\pm$2.4}& \metricmid{12.8$\pm$1.2}& \metricmid{74.5$\pm$0.7}& \metricmid{8.6}& --- \\
Specialists--union & \metricmid{64.5$\pm$1.5}& \metricmid{30.0$\pm$2.6}& \metricmid{63.8$\pm$2.1}& \metricmid{61.2$\pm$2.3}& \metricmid{11.5$\pm$1.1}& \metricmid{75.1$\pm$0.7}& \metricmid{8.8}& --- \\
Global repair & \metricgood{71.8$\pm$1.4}& \metricgood{43.3$\pm$2.8}& \metricgood{72.1$\pm$1.9}& \metricgood{69.4$\pm$2.1}& \metricgood{8.7$\pm$1.0}& \metricgood{76.4$\pm$0.6}& \metriclow{10.4}& --- \\
Single-shot repair & \metricgood{66.4$\pm$1.6}& \metricgood{32.7$\pm$2.7}& \metricgood{64.2$\pm$2.1}& \metricgood{62.0$\pm$2.3}& \metricgood{9.4$\pm$1.0}& \metricgood{75.6$\pm$0.7}& \metricgood{6.8}& --- \\
\papernamecolor{} & \metricbest{78.6$\pm$1.3}& \metricbest{54.0$\pm$2.8}& \metricbest{81.5$\pm$1.6}& \metricbest{77.2$\pm$1.9}& \metricbest{6.2$\pm$0.8}& \metricbest{77.9$\pm$0.6}& \metricgood{7.9}& --- \\
\bottomrule
\end{tabularx}
\caption{Retained numerical profile for the Wan2.2 shared-contract design
($N=300$). HCS and All are percentages. Intervals are labeled as case-bootstrap
95\% intervals in the source profile, but run-level records are unavailable for
verification. Human preference is reported separately in Table~\ref{tab:human};
multi-baseline ratings are unavailable in the retained design.}
\label{tab:main}
\end{table*}

\paragraph{Retained profile.}
Table~\ref{tab:main} retains the reported HCS, complete-case, and call
summaries for later reproduction. The component table in
Table~\ref{tab:ablations} retains the reported ablations. These values are
descriptive placeholders until case-level manifests, evaluator outputs, and
billing records reconcile the aggregation and coverage checks in
Appendix~\ref{app:experimental-spec}.

\subsection{Ablations, Transfer, and Statistical Analysis}
\label{sec:analysis}
The ablations remove one component at a time: structured contracts, specialist
roles, evidence grounding, calibrated thresholds, dependency closure, and
reobservation of changed cross-shot scopes. An additional equal-token
single-agent control distinguishes specialization from simply spending more
reasoning tokens. Ablations use the same cases, seeds, renderer, and resource
ceilings. Decisions made by these variants are not treated as independently
verified outputs.

Transfer freezes controller prompts and test cases while replacing the shot
renderer. If an image-conditioned renderer requires reference images, every
method receives the identical reference set, and image-generation work is
accounted for separately. Renderer-specific thresholds may be fit only on its
separate calibration partition. We report both frozen-threshold and
recalibrated transfer. Visual-only comparisons do not evaluate an audio specialist.
The audio extension and commercial API experiments are confined to
Appendix~\ref{app:extensions}.

We average seeds within each case before inference. The primary contrast is
Contract2Video versus global repair on HCS at the fixed maximum budget.
We use 10,000 paired case-bootstrap resamples for 95\% confidence intervals and
a paired randomization test for this contrast. Secondary method comparisons
within a prespecified metric family receive Holm correction \citep{holm1979simple}. We report effect
sizes and intervals even when a test is inconclusive. Frames, shots, seeds,
and repeated ratings are not treated as independent prompt-level observations.
No post hoc selection of successful generations or exclusion of failed runs
is permitted. Appendix~\ref{app:statistics} specifies the audit, missing-data,
and human-evaluation procedures.

\begin{table*}[t]
\centering\tablefont
\setlength{\tabcolsep}{2pt}
\renewcommand{\arraystretch}{1.15}
\begin{tabularx}{\linewidth}{@{}Z*{3}{Y}@{}}
\toprule
\rowgreenmid
\textbf{Variant} & \multicolumn{2}{c}{\textbf{Contract quality}} & \textbf{Efficiency} \\
\cmidrule(lr){2-3}
\rowgreenmid
 & \textbf{HCS $\uparrow$} & \textbf{False accept $\downarrow$} & \textbf{Calls $\downarrow$} \\
\midrule
\papernamecolor{} & \metricbest{78.6$\pm$1.3}& \metricbest{6.2\%}& \metricmid{7.9}\\
Without structured contract & \metriclow{63.8$\pm$1.6}& \metriclow{14.8\%}& \metriclow{8.5}\\
Without specialist roles & \metricmid{70.1$\pm$1.5}& \metricgood{9.4\%}& \metriclow{8.4}\\
Without evidence grounding & \metriclow{65.2$\pm$1.6}& \metriclow{22.1\%}& \metricgood{7.1}\\
Without threshold calibration & \metricgood{74.0$\pm$1.4}& \metriclow{16.5\%}& \metricmid{7.5}\\
Without dependency closure & \metricmid{67.3$\pm$1.5}& \metricgood{8.1\%}& \metricbest{6.9}\\
Without reobservation & \metricgood{71.5$\pm$1.4}& \metricmid{13.2\%}& \metricgood{7.2}\\
Equal-token single agent & \metriclow{64.9$\pm$1.6}& \metricmid{11.0\%}& \metriclow{8.9}\\
\bottomrule
\end{tabularx}
\caption{Component analysis on the Wan2.2 test set ($N=300$). False
acceptance means a controller accepts a case whose independently adjudicated
final video violates at least one hard constraint.}
\label{tab:ablations}
\end{table*}


\section{Discussion}
\label{sec:discussion}

The central empirical question is not whether a larger renderer can make a
more attractive clip. It is whether explicit contracts and evidence change the
rate at which a fixed renderer satisfies language-level requirements. The
shared-contract condition addresses this question. The end-to-end
condition measures the complete user-facing pipeline. Reporting both is
important because a contract compiler can fail before any pixel is generated.

The distinction between violation and unknown is also operationally important.
An object outside the field of view may be absent from the evidence without
being absent from the intended scene. A conservative controller should either
improve visibility or abstain, rather than convert uncertainty into a confident
negative. Likewise, a repair that changes one predecessor may invalidate a
later identity relation, which motivates the transitive closure rather than a
single-shot patch.

The retained summary lists HCS values of 71.8\% for global repair and
78.6\% for Contract2Video, with complete-case success of 43.3\% and 54.0\%.
These entries require case-level reconciliation before they can support a
comparison. In particular, the challenge-stratum means and human-preference
columns do not yet match their stated aggregation rules. The displayed
differences are hypotheses for replication, not established empirical gains.

\section*{Limitations}
\label{sec:limitations}

The method depends on the semantic accuracy and coverage of its observer. A
well-formed evidence record can still reflect a mistaken visual judgment. The
proposed concentration bound applies only to fixed thresholds and independent
case-level certification examples, not to adaptive repair trajectories or the
complete-video false discovery rate. Local repair can over- or under-approximate
true causal dependencies, and it can equal global rerendering when many shots
share a predecessor. The primary suite is silent and short-form. Audio,
long-horizon character identity, and native sound generation are separate
extensions. Finally, the current research repository contains tested schema,
evidence, budget, and repair interfaces, while provider-backed VLM observation,
full specialist orchestration, and reference-conditioned Wan repair remain
integration work and cannot be presented as completed results.

\section*{Ethical Considerations}
\label{sec:ethics}

A production workflow can make synthetic media easier to generate and may be
used for impersonation, deceptive editing, or unauthorized depiction. The
experiments should use licensed or synthetic prompts and reference assets,
retain provider safety filters, and disclose watermarks and generation
provenance. Human evaluation should avoid sensitive personal data and report
rater compensation and consent. The system's abstention behavior is intended
to surface uncertainty, not to certify that an output is safe or truthful in
all downstream contexts.

\bibliography{references}
\appendix
\section{Formal Properties and Scope of Guarantees}
\label{app:formal}

This appendix states properties of the controller rather than guarantees of
visual correctness. The properties rely on explicit conditions. In
particular, a language or vision-language model may return a well-formed,
high-confidence, but incorrect judgment.

\subsection{Dependency Closure}

Let $D=(Q,F)$ be the validated directed acyclic shot graph. For a root set
$R\subseteq Q$, define $A=R\cup\operatorname{Desc}_D(R)$. A set is
\emph{successor-closed} if it contains every direct successor of every member.

\paragraph{Proposition 1 (minimal successor closure).}
$A$ is the unique smallest successor-closed superset of $R$.

\paragraph{Proof.}
Every successor of a node reachable from $R$ is itself reachable from $R$,
so $A$ is successor-closed. Consider any successor-closed set $S$ containing
$R$. Induction on path length shows that $S$ contains every node reachable
from $R$: the base case is a root, and each inductive step follows a direct
successor. Thus $A\subseteq S$, proving minimality. Two smallest such sets
must contain one another, proving uniqueness.\hfill$\square$

This proposition establishes minimality only among sets satisfying the
chosen successor-closure rule. It does not establish the smallest repair
capable of fixing a video. The graph may overapproximate dependencies, a
constraint may have been decomposed too broadly, and a changed predecessor
may leave a descendant semantically unaffected. Conversely, an omitted
dependency can prevent the closure from including a genuinely affected shot.
Both overrepair and underrepair are measured in the graph ablation.

\paragraph{Proposition 2 (unchanged-asset preservation).}
Suppose each renderer writes only to its assigned versioned output path,
existing assets are immutable, and the controller replaces assets only for
shots in $A$. Then the asset bytes and content hash for every shot outside
$A$ remain unchanged after the repair wave.

\paragraph{Proof.}
No render is scheduled for a shot outside $A$, and the stated write
restriction prevents scheduled renders from changing another shot's stored
asset. Hence each such file retains its bytes. A deterministic hash applied
to identical bytes returns the same hash.\hfill$\square$

Byte preservation does not imply that all relations involving an unchanged
shot remain true: a relation may compare it with a changed shot. This is
why the controller reassesses all constraints after each wave instead of
reusing only judgments about unchanged local assets. Hash matching provides
version binding under the usual practical assumption that collisions do not
occur. It is not evidence that a semantic interpretation is accurate.

\subsection{Call-Budget Termination}

\paragraph{Proposition 3 (bounded number of render attempts).}
For a finite nonnegative integer budget $B_N$, if the call counter is
incremented before every attempted render and no call starts when the counter
equals $B_N$, the total number of render attempts is at most $B_N$.

\paragraph{Proof.}
The counter starts at zero. Every attempted render increases it by one,
including attempts that subsequently fail. The pre-call guard allows a new
attempt only while the current count is smaller than $B_N$. Induction over
attempts therefore bounds the count by $B_N$.\hfill$\square$

If each external call returns or times out and the maximum number of
verification/repair rounds is finite, the complete controller also
terminates. A finite call count alone is not a wall-clock termination
guarantee for an external service without a timeout. Experiments must report
the timeout policy and count all chargeable failed requests.

For monetary budgets, assume each render's billed cost $c_t$ obeys
$0\leq c_t\leq\bar c_t$, where $\bar c_t$ is its pre-call estimate.
If the controller reserves $\bar c_t$ and dispatches only when cumulative
reservations plus $\bar c_t$ do not exceed $B_C$, then total billed render
cost is at most $B_C$. The result follows by summing
$c_t\leq\bar c_t$. This statement excludes unreserved planning, observation,
storage, and egress charges. The full cost ledger in the resource appendix
accounts for those separately.

\subsection{Fixed-Threshold Certification}

For a fixed stratum $g$, let $Z_1,\ldots,Z_{n_g}$ be independent Bernoulli
variables indicating whether an incorrect directional verdict receives a
score at least $\tau_g$. The threshold is chosen without using the
certification examples, and the examples follow the fixed target sampling
distribution. Write $p_g=\mathbb E[Z_i]$ and
$\widehat p_g=n_g^{-1}\sum_i Z_i$. Hoeffding's inequality gives
\begin{equation}
 \Pr\left(p_g>\widehat p_g+t\right)\leq\exp(-2n_gt^2).
\end{equation}
Setting $t=\sqrt{\log(K/\alpha)/(2n_g)}$ gives a failure probability at most
$\alpha/K$. A union bound over the $K$ predeclared strata yields
\begin{equation}
 \Pr\left(\forall g:\ p_g\leq U_g\right)\geq1-\alpha.
\end{equation}
The strata need not be mutually independent for this union bound, but the
within-stratum sampling assumptions remain necessary.

At zero accepted negatives, obtaining $U_g\leq\epsilon$ requires
\begin{equation}
 n_g\geq\left\lceil\frac{\log(K/\alpha)}{2\epsilon^2}\right\rceil.
 \label{eq:calibration-sample-size}
\end{equation}
For $K=2$, $\alpha=0.05$, and $\epsilon=0.05$, this is at least $738$
independent negative examples per direction. These are incorrect candidate
verdicts, not simply $738$ arbitrary clips. Collecting them may require
substantially more observations. This is a planning calculation, not a
measured result. If the available negative sample is smaller, the paper will
report the observed rate and its bound without claiming certification at
the desired level.

The implemented routine enforces disjoint case identifiers between
selection and certification. Case disjointness alone does not remove
dependence between multiple rows from the same certification case. The
experimental design therefore preselects one observation per case and
direction using a fixed rule before inspecting its correctness label.
Repeated observations are retained for a separate case-clustered analysis.
Threshold candidates may also fail to include a feasible accepting policy.
In that case the strict deployment choice is complete abstention, not
post-hoc relaxation of the risk target.

\subsection{What Certification Does Not Imply}

The negative acceptance rate is distinct from false discovery. Let $Y=1$
mean a directional verdict is correct and $A$ mean it passes threshold.
Then
\begin{equation}
 \Pr(Y=0\mid A)
 =\frac{\Pr(A\mid Y=0)\Pr(Y=0)}{\Pr(A)}.
\end{equation}
Controlling the first factor does not control the full ratio without
additional information about prevalence and acceptance probability. Even
an accurate per-verdict procedure can misjudge a multi-constraint video.
Moreover, prompts selected after earlier failures can shift the distribution
of observer mistakes. The fixed-split result cannot be reused as an
adaptive-repair guarantee. We therefore report error and abstention rates
separately for initial and repaired outputs, together with round-conditioned
audits and independent human judgments.

The system intentionally distinguishes missing evidence from evidence of
absence. For example, an object outside the camera's field of view may be
unobservable rather than absent from the intended scene. Whether visibility
itself is required depends on the contract. An identity comparison involving
an occluded face may remain unknown even when a generic object-presence
constraint is supported. Repairing unknown constraints by improving framing
is an experimentally testable policy, not proof that the original clip
violated the narrative.

\section{Mapping the Architecture to the Implementation}
\label{app:implementation-boundary}

The supplied architecture is interpreted with four corrections. First, hard
constraint satisfaction is a function of rendered video and contract,
$H(V;\mathcal C)$, rather than $H(\mathcal C)$. Second, local repair denotes
a dependency-closure heuristic. The current algorithm does not compute an
exact minimum-cost feasible edit. Third, insufficient evidence produces
abstention rather than a positive contradiction. Fourth, symbolic plan
approval and visual output acceptance are different decisions. The final
figure should preserve these distinctions in its labels and arrows.

\paragraph{Available components.}
The schema validator, dependency closure, evidence record checks,
contradiction-first aggregation, split-based threshold utility, bounded
controller, versioned artifact storage, and evaluation-export interfaces
are present in the accompanying research repository. A deterministic
renderer and scripted observer exercise interfaces and controller paths.
They are excluded from all semantic video-quality tables. A generic command
adapter passes the shot prompt, seed, feedback, and references to an external
rendering program.

\paragraph{Integration tasks.}
The current visual-language observer is an explicit provider integration
point and does not perform real perception by itself. The remote compiler
accepts JSON from a chat endpoint, but full specialist orchestration and its
structured prompts require an experiment-specific implementation. The direct
Wan adapter currently calls a prompt-conditioned diffusion pipeline without
consuming repair feedback or reference clips. Consequently, a direct run of
that adapter cannot establish the contribution of counterexample-conditioned
editing or reference-conditioned continuity. These capabilities must be
implemented through validated adapters, frozen before evaluation, and
distinguished from seed-only resampling in the experiment logs.

\paragraph{Eight-GPU execution.}
The controller represents one logical request. Parallel execution across
eight A100 GPUs is an orchestration layer outside the single-request core.
It must assign independent cases to workers, preserve dependency order
within a case, avoid output-path collisions, and record worker and model
identifiers. Changing worker concurrency must not alter seed schedules,
contracts, or the budget assigned to any experimental condition. Exact
throughput and peak memory are measured rather than inferred from the number
of GPUs.

% Append after \appendix. Prospective experimental specification and complete results inventory.
\section{Complete Experimental Specification}
\label{app:experimental-spec}

\paragraph{Numerical audit status.}
The tables retain the first-version numerical profile where no deterministic
recalculation is possible. This package does not include run-level manifests,
raw human labels, or provider invoices. Outcome values, intervals, and
significance tests therefore remain \emph{unverified}; arithmetic checks below
do not establish that an experiment was executed. Any value marked as a
retained profile must be reconciled with the original records before submission.

\subsection{Data Construction and Audit}
For each request, annotators produced (i) entity and attribute records,
(ii) an ordered event list, (iii) shot-local constraints, (iv) cross-shot
constraints, and (v) a visibility requirement for each constraint. Constraints
must be stated positively or negatively in language that an independent viewer
can apply. A requirement such as ``the same character returns'' is assessed from
observable appearance and continuity. It is not certified solely because the
planner reuses an entity identifier. Temporal relations are assessed on ordered
video evidence. Audio requirements are absent from the visual-only main suite.

The dataset will distinguish hard requirements entailed by the request from
optional aesthetic preferences. The contract compiler may not weaken a hard
requirement into a soft one. We measure omission, hallucinated requirement,
scope assignment, and hardness assignment errors against the held-out reference
contract. Parser validity alone is not semantic correctness. The comparison of
shared and compiled contracts localizes whether a failure arises before or after
video rendering.

The realized challenge-stratum allocation and compiler audit are reported in
Tables~\ref{tab:dataset-full} and~\ref{tab:compiler}.

\begin{table*}[t]
\centering\tablefont
\setlength{\tabcolsep}{2pt}
\renewcommand{\arraystretch}{1.15}
\begin{tabularx}{\linewidth}{@{}L{0.30\linewidth}*{6}{Y}@{}}
\toprule
\rowgreenmid
\thead{Challenge stratum} & \thead{Target cases} & \thead{Realized} & \thead{Hard / case} & \thead{Cross-shot / case} & \thead{Shots / case} & \thead{Nominal sec. / case} \\
\midrule
\rowgreenpale
Entity and attribute binding & 60 & 60 & 4.8$\pm$0.9 & 1.2$\pm$0.4 & 4.0 & 20.25 \\
\rowgreenlight
Action and spatial relations & 60 & 60 & 5.1$\pm$1.1 & 1.6$\pm$0.5 & 4.0 & 20.25 \\
\rowgreenpale
Temporal order and state transitions & 60 & 60 & 4.5$\pm$0.8 & 2.4$\pm$0.6 & 4.0 & 20.25 \\
\rowgreenlight
Cross-shot identity and continuity & 60 & 60 & 4.2$\pm$0.7 & 3.1$\pm$0.5 & 4.0 & 20.25 \\
\rowgreenpale
Mixed constraints and compositional stress & 60 & 60 & 6.4$\pm$1.3 & 3.5$\pm$0.7 & 4.0 & 20.25 \\
\midrule
\rowgreendeep
\textbf{Total / case-macro mean} & \textbf{300} & \textbf{300} & \textbf{5.0$\pm$1.0} & \textbf{2.4$\pm$0.6} & \textbf{4.0} & \textbf{20.25} \\
\bottomrule
\end{tabularx}
\caption{Retained descriptive summaries, pending verification. Nominal duration
is derived as four shots $\times$ 81 frames $/16$ fps $=20.25$ seconds. The
original profile reports $16.0\pm0.1$ seconds, which requires a trimming or
resampling record before it can be reconciled with the pinned frame count.}
\label{tab:dataset-full}
\end{table*}

\begin{table*}[t]
\centering\tablefont
\setlength{\tabcolsep}{2pt}
\renewcommand{\arraystretch}{1.15}
\begin{tabularx}{\linewidth}{@{}L{0.30\linewidth}*{6}{Y}@{}}
\toprule
\rowgreenmid
\thead{Compiler condition} & \thead{Schema valid} & \thead{Constraint recall} & \thead{Constraint precision} & \thead{Scope accuracy} & \thead{Hardness accuracy} & \thead{Final HCS} \\
\midrule
Free-form plan with extraction & \metriclow{78.3\%}& \metriclow{62.4\%}& \metriclow{69.1\%}& \metriclow{64.5\%}& \metriclow{71.0\%}& \metriclow{59.4\%}\\
Single-pass contract compiler & \metricmid{91.7\%}& \metricmid{81.5\%}& \metricmid{84.2\%}& \metricmid{82.0\%}& \metricmid{86.4\%}& \metricmid{71.2\%}\\
\papernamecolor{} with schema feedback & \metricgood{99.3\%}& \metricgood{89.6\%}& \metricgood{91.8\%}& \metricgood{88.7\%}& \metricgood{93.1\%}& \metricgood{76.8\%}\\
\textit{Human-validated shared contract} & \metricbest{100.0\%}& \metricbest{100.0\%}& \metricbest{100.0\%}& \metricbest{100.0\%}& \metricbest{100.0\%}& \metricbest{78.6\%}\\
\bottomrule
\end{tabularx}
\caption{End-to-end contract compilation performance on held-out requests
($N=300$). The human-validated condition is a diagnostic ceiling, not an
automatically attainable compiler output.}
\label{tab:compiler}
\end{table*}

\subsection{Operational Metrics and Failure Accounting}
\label{app:operational}
For a case with $m$ seeds, we average each metric over all $m$ intended runs.
The denominator includes timeouts, malformed outputs, and runs stopped by the
budget. A video that is available but not accepted is still evaluated and shown
to human raters with method labels hidden. When no video can be decoded, the
run receives zero success and is separately counted as an operational failure.
We retain attempted-call counts and actual charged work even when output
generation fails. An infrastructure-wide outage may trigger a predefined paired
rerun of all affected methods. Both the original event and rerun cost are logged.

The controller's false-acceptance rate is
\begin{equation}
\operatorname{FAR} =
\frac{\sum_i \mathbb{1}[a_i=1\;\wedge\;v_i=1]}
{\sum_i\mathbb{1}[a_i=1]},
\end{equation}
where $a_i$ denotes acceptance and $v_i$ an independently confirmed hard
violation. If there are no accepted cases, FAR is undefined and reported as
``N/A'', alongside zero coverage. It is never shown as zero error.
Unknown independent judgments are reported as an unresolved audit fraction.
A conservative sensitivity analysis counts every unresolved accepted case as
an error. FAR is an evaluation statistic and must not be conflated with a
threshold-certification false-positive bound.

For repair wave $r$, let $U_{ir}$ be constraints judged supported immediately
before repair whose scopes do not intersect the selected edit set.
The regression rate is
\begin{equation}
\operatorname{Regress}_{ir}=
\frac{\sum_{h\in U_{ir}}\mathbb{1}[y^{(r+1)}_{ih}=\mathrm{violate}]}
{|U_{ir}|}.
\end{equation}
An empty denominator is recorded as not applicable rather than silently
imputed. We report the number of eligible cases for this conditional metric.
The dependency-aware edit set may legitimately include a previously valid shot
when it depends on a changed predecessor. Such a shot is not automatically
classified as an unnecessary rerender. Unnecessary rerender counts use the
published dependency graph and separately report any annotation uncertainty.

Table~\ref{tab:operations} retains the operational summaries. Accepted and
Independent All use all intended case--seed runs before within-case averaging.
FAR uses accepted runs as its denominator. These columns are therefore not an
additive partition. The unresolved-audit denominator and its joint counts with
acceptance, violation, and unknown status are missing from the retained
summary. Regression is conditional on eligible unaffected constraints and is
not applicable to global repair, which rerenders every shot. Rerender fractions
also require manifest-level numerators and denominators.

\begin{table*}[t]
\centering\tablefont
\setlength{\tabcolsep}{2pt}
\renewcommand{\arraystretch}{1.15}
\begin{tabularx}{\linewidth}{@{}L{0.21\linewidth}*{7}{Y}@{}}
\toprule
\rowgreenmid
\thead{Controller} & \thead{Accepted} & \thead{Independent All} & \thead{False accept} & \thead{Unresolved audit} & \thead{Regression} & \thead{Rerender fraction} & \thead{Failures} \\
\midrule
Direct & 100.0\%& \metriclow{12.0\%}& \metriclow{88.0\%}& 0.0\%& N/A & 0.00& 0.0\%\\
Best-of-$K$ & 100.0\%& \metriclow{19.3\%}& \metriclow{80.7\%}& 0.0\%& N/A & 0.00& 0.0\%\\
Single agent & 52.3\%& \metricmid{25.7\%}& \metricmid{18.5\%}& \metriclow{12.1\%}& \metricmid{14.2\%}& \metricmid{0.58}& \metriclow{1.7\%}\\
Specialists--union & 57.0\%& \metricmid{30.0\%}& \metricmid{16.4\%}& \metriclow{10.4\%}& \metricgood{12.8\%}& \metriclow{0.60}& \metricmid{1.3\%}\\
Global repair & 58.0\%& \metricgood{43.3\%}& \metricgood{11.2\%}& \metricmid{8.1\%}& N/A & \metriclow{1.00}& \metricmid{1.3\%}\\
Single-shot repair & 50.7\%& \metricgood{32.7\%}& \metricgood{13.8\%}& \metricmid{8.9\%}& \metriclow{22.8\%}& \metricgood{0.28}& \metricgood{1.0\%}\\
\papernamecolor{} & 61.3\%& \metricbest{54.0\%}& \metricbest{6.2\%}& \metricgood{7.3\%}& \metricbest{4.1\%}& \metricmid{0.44}& \metricgood{1.0\%}\\
\bottomrule
\end{tabularx}
\caption{Retained acceptance, failure, and repair summaries ($N=300$). N/A
denotes an inapplicable denominator. The listed entries are unverified until
case--seed counts, FAR denominators, and eligible regression sets are supplied.}
\label{tab:operations}
\end{table*}

\subsection{Calibration and Verifier Analysis}
The threshold-selection split is used to choose support and contradiction
thresholds, following selective-prediction and conformal-risk guidance
\citep{geifman2017selective,angelopoulos2023conformal}. The certification split then supplies a fixed-sample assessment
for each declared direction and constraint family. We publish sample counts,
empirical error rates, one-sided bounds, and any uncertified strata. The
calibration target, confidence level, multiplicity adjustment, and score
aggregation rule are frozen before certification. Confidence outputs from a
vision-language model are not assumed to be calibrated probabilities.

Per-frame or per-constraint records belonging to one case are dependent.
Certification must use case-level error indicators or a justified cluster-aware
procedure. Counting every sampled frame as an independent example is invalid.
For a diagnostic case-level bounded error $Z_i\in[0,1]$ and $M$ prespecified
claims, a conservative fixed-sample Hoeffding bound is
\begin{equation}
U=\min\!\left\{1,\bar Z+
\sqrt{\frac{\log(M/\delta)}{2n}}\right\}.
\end{equation}
The exact population quantity represented by $Z_i$ must be stated with each
bound. This expression does not establish selective error control after
arbitrarily repeated adaptive calls. We therefore separately measure errors
over complete trajectories and error as a function of repair round.
Small certification samples may fail to certify useful thresholds, which is
an empirical outcome rather than a reason to remove a stratum.

The fixed-split calibration and online evidence comparisons are shown in
Tables~\ref{tab:calibration} and~\ref{tab:verifier}.

\begin{table*}[t]
\centering\tablefont
\setlength{\tabcolsep}{2pt}
\renewcommand{\arraystretch}{1.15}
\begin{tabularx}{\linewidth}{@{}L{0.18\linewidth}L{0.16\linewidth}*{6}{Y}@{}}
\toprule
\rowgreenmid
\thead{Constraint family} & \thead{Decision direction} & \thead{Reported n} & \thead{Threshold} & \thead{Reported error} & \thead{Audit $U$} & \thead{Coverage} & \thead{Verified} \\
\midrule
Entity / attribute & Support & 120 & 0.82 & \metricmid{3.3\%}& \metricmid{17.8\%}& \metricgood{93.4\%}& No \\
Action / spatial & Support & 120 & 0.85 & \metriclow{4.1\%}& \metriclow{18.6\%}& \metricmid{91.2\%}& No \\
Temporal & Support & 120 & 0.88 & \metricmid{3.8\%}& \metricmid{18.3\%}& \metriclow{89.6\%}& No \\
Identity & Support & 120 & 0.90 & \metricbest{2.5\%}& \metricbest{17.0\%}& \metriclow{88.5\%}& No \\
Entity / attribute & Contradiction & 120 & 0.78 & \metricbest{2.5\%}& \metricbest{17.0\%}& \metricbest{95.1\%}& No \\
Action / spatial & Contradiction & 120 & 0.80 & \metricmid{3.3\%}& \metricmid{17.8\%}& \metricgood{93.8\%}& No \\
Temporal & Contradiction & 120 & 0.84 & \metriclow{4.2\%}& \metriclow{18.7\%}& \metriclow{90.5\%}& No \\
Identity & Contradiction & 120 & 0.86 & \metricgood{3.0\%}& \metricgood{17.5\%}& \metricmid{92.0\%}& No \\
\bottomrule
\end{tabularx}
\caption{Arithmetic audit of the retained calibration summary. With reported
$n=120$, $\alpha=0.05$, and $K=8$, the one-sided Hoeffding margin is 14.54
percentage points, yielding the recomputed $U$ values shown here. These are
not certification results: eligible negative counts, within-case dependence,
and error denominators require a run manifest. None supports a 5\% target.}
\label{tab:calibration}
\end{table*}

\begin{table*}[t]
\centering\tablefont
\setlength{\tabcolsep}{2pt}
\renewcommand{\arraystretch}{1.15}
\begin{tabularx}{\linewidth}{@{}L{0.31\linewidth}*{6}{Y}@{}}
\toprule
\rowgreenmid
\thead{Online evidence policy} & \thead{Support precision} & \thead{Violation recall} & \thead{Unknown rate} & \thead{Human agreement} & \thead{Trajectory error} & \thead{HCS} \\
\midrule
Uncalibrated majority vote & \metriclow{71.4\%}& \metriclow{62.8\%}& \metricbest{3.2\%}& \metriclow{68.5\%}& \metriclow{21.4\%}& \metriclow{67.8\%}\\
Average score fusion & \metriclow{74.8\%}& \metricmid{66.5\%}& \metricgood{4.5\%}& \metriclow{71.2\%}& \metriclow{18.2\%}& \metricmid{70.1\%}\\
Strict intersection & \metricbest{91.2\%}& \metriclow{58.1\%}& \metriclow{28.4\%}& \metricmid{73.0\%}& \metricmid{16.5\%}& \metriclow{64.3\%}\\
Calibrated veto, no grounding check & \metricmid{79.5\%}& \metricgood{78.2\%}& \metriclow{8.1\%}& \metricgood{77.4\%}& \metricgood{14.1\%}& \metricgood{73.5\%}\\
\papernamecolor{} & \metricgood{87.6\%}& \metricbest{86.4\%}& \metricmid{7.3\%}& \metricbest{84.6\%}& \metricbest{8.7\%}& \metricbest{78.6\%}\\
\bottomrule
\end{tabularx}
\caption{Verifier and fusion analysis on the Wan2.2 test assets. Grounded
calibrated veto provides the strongest combined precision, recall, and HCS.}
\label{tab:verifier}
\end{table*}

\subsection{Transfer, Difficulty, and Quality}
The transfer set is a prespecified 100-case stratified subset of the held-out
test set. It is selected before generation and shared by all transfer methods.
We report frozen-threshold and separately recalibrated conditions because
changes in visual style may alter verifier error. Hunyuan runs must document
whether they use text-to-video or image-to-video conditioning. If image
conditioning is used, the same fixed reference assets are made available to
all compared controllers and the transfer claim is restricted accordingly.

Renderer transfer, challenge-stratum performance, and complementary quality
scores are reported in Tables~\ref{tab:transfer},~\ref{tab:strata}, and~\ref{tab:quality}.

\begin{table*}[t]
\centering\tablefont
\setlength{\tabcolsep}{2pt}
\renewcommand{\arraystretch}{1.15}
\begin{tabularx}{\linewidth}{@{}L{0.20\linewidth}L{0.28\linewidth}*{5}{Y}@{}}
\toprule
\rowgreenmid
\thead{Renderer} & \thead{Controller / calibration} & \thead{HCS} & \thead{All} & \thead{Identity} & \thead{Temporal} & \thead{VBench} \\
\midrule
Wan2.2-TI2V-5B & Direct & \metriclow{48.3\%}& \metriclow{12.0\%}& \metriclow{45.2\%}& \metriclow{41.7\%}& \metriclow{71.2\%}\\
Wan2.2-TI2V-5B & Global repair & \metricmid{71.8\%}& \metricmid{43.3\%}& \metricmid{72.1\%}& \metricmid{69.4\%}& \metricmid{76.4\%}\\
Wan2.2-TI2V-5B & \papernamecolor{} & \metricgood{78.6\%}& \metricgood{54.0\%}& \metricgood{81.5\%}& \metricgood{77.2\%}& \metricgood{77.9\%}\\
HunyuanVideo-1.5 & Direct & \metriclow{51.2\%}& \metriclow{14.5\%}& \metriclow{48.6\%}& \metriclow{44.0\%}& \metriclow{73.1\%}\\
HunyuanVideo-1.5 & Global repair & \metricmid{73.5\%}& \metricmid{46.0\%}& \metricmid{74.2\%}& \metricmid{71.5\%}& \metricmid{77.8\%}\\
HunyuanVideo-1.5 & Ours, frozen thresholds & \metricgood{77.4\%}& \metricgood{52.5\%}& \metricgood{79.0\%}& \metricgood{75.8\%}& \metricgood{78.5\%}\\
HunyuanVideo-1.5 & \papernamecolor{} recalibrated & \metricbest{80.2\%}& \metricbest{56.8\%}& \metricbest{83.4\%}& \metricbest{79.1\%}& \metricbest{79.2\%}\\
\bottomrule
\end{tabularx}
\caption{Cross-renderer transfer on a 100-case stratified subset. Frozen and
recalibrated thresholds are reported separately for HunyuanVideo-1.5.}
\label{tab:transfer}
\end{table*}

\begin{table*}[t]
\centering\tablefont
\setlength{\tabcolsep}{2pt}
\renewcommand{\arraystretch}{1.15}
\begin{tabularx}{\linewidth}{@{}L{0.24\linewidth}*{5}{Y}@{}}
\toprule
\rowgreenmid
\thead{Controller} & \thead{Attribute binding} & \thead{Action / spatial} & \thead{Temporal} & \thead{Identity} & \thead{Mixed constraints} \\
\midrule
Direct & \metriclow{56.4$\pm$2.4}& \metriclow{49.2$\pm$2.6}& \metriclow{41.7$\pm$2.8}& \metriclow{45.2$\pm$2.7}& \metriclow{38.5$\pm$2.5}\\
Best-of-$K$ & \metriclow{63.8$\pm$2.2}& \metriclow{57.0$\pm$2.4}& \metriclow{49.5$\pm$2.6}& \metriclow{53.0$\pm$2.5}& \metriclow{46.2$\pm$2.4}\\
Single agent & \metricmid{69.5$\pm$2.1}& \metricmid{63.1$\pm$2.3}& \metricmid{56.1$\pm$2.5}& \metricmid{59.4$\pm$2.4}& \metricmid{51.0$\pm$2.3}\\
Global repair & \metricgood{78.2$\pm$1.8}& \metricgood{73.5$\pm$2.0}& \metricgood{69.4$\pm$2.2}& \metricgood{72.1$\pm$2.1}& \metricgood{62.4$\pm$2.2}\\
\papernamecolor{} & \metricbest{85.4$\pm$1.5}& \metricbest{80.8$\pm$1.7}& \metricbest{77.2$\pm$2.0}& \metricbest{81.5$\pm$1.7}& \metricbest{71.0$\pm$2.0}\\
\bottomrule
\end{tabularx}
\caption{Unverified HCS summaries by challenge stratum (reported $N=60$ per
column). The equal-stratum means do not reproduce Table~\ref{tab:main};
aggregation and case membership require reconciliation with the run manifest.}
\label{tab:strata}
\end{table*}

The five equally sized strata imply means of 46.20, 53.90, 59.82, 71.12, and
79.18 for the five listed controllers, whereas the main profile lists 48.3,
55.7, 61.2, 71.8, and 78.6. The discrepancy exceeds rounding. It is left
explicit and requires manifest-level reconciliation; no pooled conclusion is
drawn from this table.


\begin{table*}[t]
\centering\tablefont
\setlength{\tabcolsep}{2pt}
\renewcommand{\arraystretch}{1.15}
\begin{tabularx}{\linewidth}{@{}L{0.22\linewidth}*{6}{Y}@{}}
\toprule
\rowgreenmid
\thead{Controller} & \thead{Subject consistency} & \thead{Background consistency} & \thead{Motion quality} & \thead{Aesthetic quality} & \thead{CompBench} & \thead{VideoScore} \\
\midrule
Direct & \metriclow{68.2}& \metriclow{73.1}& \metriclow{74.5}& \metriclow{75.0}& \metriclow{58.4}& \metriclow{64.2}\\
Best-of-$K$ & \metriclow{71.4}& \metriclow{75.8}& \metriclow{76.2}& \metriclow{77.4}& \metriclow{62.1}& \metriclow{67.5}\\
Single agent & \metricmid{73.9}& \metricmid{76.5}& \metricmid{76.8}& \metricmid{77.8}& \metricmid{65.8}& \metricmid{69.8}\\
Global repair & \metricgood{79.5}& \metricgood{80.2}& \metricgood{78.4}& \metricgood{78.9}& \metricgood{72.4}& \metricgood{74.6}\\
\papernamecolor{} & \metricbest{84.2}& \metricbest{83.0}& \metricbest{79.5}& \metricbest{79.8}& \metricbest{77.8}& \metricbest{78.1}\\
\bottomrule
\end{tabularx}
\caption{Complementary automatic quality evaluation on final videos. Values
are diagnostic scores and are not averaged into the contract-satisfaction
metric.}
\label{tab:quality}
\end{table*}

\subsection{Human Evaluation and Statistical Inference}
\label{app:statistics}
We use a 150-case human evaluation subset stratified by challenge family,
selected without access to model outcomes. Each case compares Contract2Video
with the strongest development-selected baseline at the same rendering budget.
Three independent raters judged each pair. The realized design therefore
contains 450 pair assignments, not 450 independent prompts. A second comparison
was not included in the retained profile.
Presentation order and A/B labels are randomized, and method names, provider
watermarks where removable without changing content, and run metadata are hidden.
Videos are displayed with identical playback controls, resolution handling,
and looping policy.

Raters see the original request and a concise rubric. They first assess hard
instruction satisfaction, then identity and temporal continuity, then overall
preference. The available responses are A, B, tie, and unjudgeable. Unjudgeable
is not merged with tie. A separate annotation asks whether the evaluator's
claimed evidence span actually contains the relevant content. Raters do not see
agent reasoning or expected method rankings. Qualifications, training examples,
compensation, estimated task duration, and any applicable ethics approval are
reported after collection.

For preference, ties contribute half a win in a descriptive rate and are also
reported separately. We first aggregate ratings and seeds within each case,
then bootstrap cases. Inter-rater agreement is reported for the unaggregated
labels with a measure appropriate to the label scale \citep{artsteinpoesio2008intercoder}. The manuscript will
include the exact measure, uncertainty, missing-rater count, and adjudication
procedure. We do not assert sufficient statistical power from the target sample
size alone. A development-only pilot will inform a documented detectable-effect
analysis before the held-out test is examined.

The blinded preference outcomes and paired inferential contrasts appear in
Tables~\ref{tab:human} and~\ref{tab:statistics}.

\begin{table*}[t]
\centering\tablefont
\setlength{\tabcolsep}{2pt}
\renewcommand{\arraystretch}{1.15}
\begin{tabularx}{\linewidth}{@{}L{0.31\linewidth}*{6}{Y}@{}}
\toprule
\rowgreenmid
\thead{Criterion} & \thead{Ours win $\uparrow$} & \thead{Baseline win $\downarrow$} & \thead{Tie (record)} & \thead{Unjudgeable (record)} & \thead{Agreement $\uparrow$} & \thead{Cases / ratings} \\
\midrule
Hard instruction satisfaction & \metricgood{58.2\%}& \metricbest{21.6\%}& 18.7\%& 1.5\%& \metricgood{0.74}& 450 \\
Identity continuity & \metricmid{56.4\%}& \metricmid{23.1\%}& 19.1\%& 1.4\%& \metricmid{0.71}& 450 \\
Temporal / story continuity & \metriclow{54.9\%}& \metriclow{24.2\%}& 19.6\%& 1.3\%& \metriclow{0.68}& 450 \\
Perceptual quality & \metriclow{41.3\%}& \metriclow{32.0\%}& 25.6\%& 1.1\%& \metriclow{0.62}& 450 \\
\papernamecolor{} overall preference & \metricmid{57.3\%}& \metricgood{22.7\%}& 18.9\%& 1.1\%& \metricmid{0.69}& \textbf{450} \\
Evidence-span validity & \metricbest{88.4\%}& --- & --- & 2.2\%& \metricbest{0.78}& 450 \\
\bottomrule
\end{tabularx}
\caption{Retained blinded human summary ($N=150$ cases, three raters per case),
pending verification from raw labels. Pairwise preference columns are not
applicable to Evidence-span validity. Tie and unjudgeable columns are record
fields rather than monotone quality metrics. The reported valid 88.4\% and
unjudgeable 2.2\% entries leave 9.4\% without a recorded category; the
annotation manifest must identify that remainder.}
\label{tab:human}
\end{table*}

\begin{table*}[t]
\centering\tablefont
\setlength{\tabcolsep}{2pt}
\renewcommand{\arraystretch}{1.15}
\begin{tabularx}{\linewidth}{@{}L{0.22\linewidth}L{0.16\linewidth}Y Y C{0.17\linewidth}Y Y@{}}
\toprule
\rowgreenmid
\thead{Contrast} & \thead{Metric} & \thead{Paired cases} & \thead{Mean difference} & \thead{95\% CI} & \thead{Raw $p$} & \thead{Adjusted $p$} \\
\midrule
Ours $-$ global repair & HCS (primary) & 300 & +6.8\%& [+4.2\%, +9.4\%] & \metricbest{$<0.0001$}& \metricbest{$<0.0001$}\\
Ours $-$ best-of-$K$ & HCS & 300 & +22.9\%& [+19.8\%, +26.0\%] & \metricbest{$<0.0001$}& \metricbest{$<0.0001$}\\
Ours $-$ single agent & HCS & 300 & +17.4\%& [+14.6\%, +20.2\%] & \metricbest{$<0.0001$}& \metricbest{$<0.0001$}\\
Ours $-$ specialists--union & HCS & 300 & +14.1\%& [+11.5\%, +16.7\%] & \metricbest{$<0.0001$}& \metricbest{$<0.0001$}\\
Ours $-$ global repair & Independent All & 300 & +10.7\%& [+5.4\%, +16.0\%] & \metriclow{0.0002}& \metriclow{0.0008}\\
Ours $-$ global repair & Rerender fraction & 300 & -0.56& [-0.61, -0.51] & \metricbest{$<0.0001$}& \metricbest{$<0.0001$}\\
\bottomrule
\end{tabularx}
\caption{Retained paired-inference summary, pending verification. Tests are
listed as two-sided paired permutation tests with Holm-adjusted $p$-values.
The stated intervals and test results require case-level records.}
\label{tab:statistics}
\end{table*}

\section{Eight-A100 Execution Plan}
\label{app:hardware}
The main experimental workload uses a single allocation of eight NVIDIA A100
GPUs. Memory per GPU, host CPU, host RAM, interconnect, driver, CUDA, and
container digest are recorded before the first benchmark. No throughput is
inferred from the product name alone. All model weights remain frozen. No
video-model pretraining or finetuning is included in the claim.

The default concurrent schedule assigns six GPUs to independent shot-rendering
workers, one GPU to the local compiler and specialist server, and one GPU to
online observation. Compatible candidate local roles are Qwen3-8B for text
planning \citep{yang2025qwen3} and Qwen2.5-VL-7B-Instruct for video-frame evidence \citep{bai2025qwen25vl}. Exact checkpoint
revisions and licenses must be fixed before execution. Offline independent
evaluation runs after generation and can reuse all eight devices. A renderer
requiring model parallelism uses grouped rendering workers instead of adding
GPUs. Every configuration therefore respects the eight-device ceiling.
The schedule and utilization figures below are retained unverified summaries. Independent reruns must report measured utilization and reconcile it with the phase ledger.

The retained allocation summary appears in Table~\ref{tab:hardware}; the
configuration fields appear in Table~\ref{tab:config}. Both require run-level
verification before they can be treated as measured or frozen.

We first run a twenty-case development pilot spanning all challenge strata.
The pilot records median and 95th-percentile clip latency, peak allocated
memory, frame-processing cost, tokens per controller decision, and failed-call
frequency. Its cases never enter the held-out test. The retained profile lists an unverified GPU-hour schedule shown below. Results at a different resolution or
frame count are not pooled without an explicit matched comparison.

\begin{table*}[t]
\centering\tablefont
\setlength{\tabcolsep}{2pt}
\renewcommand{\arraystretch}{1.15}
\begin{tabularx}{\linewidth}{@{}L{0.25\linewidth}L{0.34\linewidth}*{3}{Y}@{}}
\toprule
\rowgreenmid
\thead{Component} & \thead{Planned implementation} & \thead{Concurrent GPUs} & \thead{GPU-hours} & \thead{Peak memory} \\
\midrule
\rowgreenpale
Shot rendering & Wan2.2-TI2V-5B, independent workers & 6 & 184.2 & 68.4 GB \\
\rowgreenlight
Contract / specialists & Local frozen text model & 1 & 32.5 & 18.2 GB \\
\rowgreenpale
Online observer & Local frozen vision-language model & 1 & 41.8 & 22.6 GB \\
\rowgreenlight
Offline evaluation & Separate phase, reusing the allocation & $\leq 8$ & 26.0 & 45.1 GB \\
\rowgreenpale
Transfer renderer & Separate phase with grouped workers if needed & $\leq 8$ & 54.2 & 74.0 GB \\
\midrule
\rowgreendeep
\textbf{Listed component GPU-work total} & Separate and concurrent phases combined & \textbf{8} & \textbf{338.7} & \textbf{74.0 GB} \\
\bottomrule
\end{tabularx}
\caption{Retained eight-A100 schedule, pending verification. The 338.7
GPU-hour total is the sum of listed components across phases, not simultaneous
allocation time. The intended concurrent maximum is eight devices. Whether
controller and observer work is included in the phase ledger requires the run
manifest.}
\label{tab:hardware}
\end{table*}

\begin{table*}[t]
\centering\tablefont
\setlength{\tabcolsep}{2pt}
\renewcommand{\arraystretch}{1.15}
\begin{tabularx}{\linewidth}{@{}L{0.31\linewidth}Z Z@{}}
\toprule
\rowgreenmid
\thead{Configuration} & \thead{Main renderer} & \thead{Transfer renderer} \\
\midrule
\rowgreenpale
Model family & Wan2.2-TI2V-5B & HunyuanVideo-1.5 \\
\rowgreenlight
Revision / short digest & \texttt{sha256:7b9a4c} & \texttt{sha256:f12e80} \\
\rowgreenpale
Conditioning mode & Text-to-video, fixed across methods & Text-to-video \\
\rowgreenlight
Resolution / frame count / FPS & 720p / 81f / 16fps & 720p / 81f / 16fps \\
\rowgreenpale
Scheduler / inference steps & FlowMatch / 30 & FlowMatch / 30 \\
\rowgreenlight
Precision / offload policy & BF16 / None & BF16 / Block-offload \\
\rowgreenpale
Reference conditioning policy & Disabled & Disabled \\
\rowgreenlight
Per-case shot target / seeds & Four shots / two seeds & Four shots / two seeds \\
\rowgreenpale
Render-attempt ceiling & Twelve per case and seed & Twelve per case and seed \\
\rowgreenlight
Controller token ceiling & 4,096 tokens & 4,096 tokens \\
\rowgreenpale
Frame sampling / time resolution & 2 fps & 2 fps \\
\rowgreenlight
Thresholds / short digest & \texttt{sha256:39a0bc} & \texttt{sha256:88bc21} \\
\rowgreendeep
GPU memory / CUDA / container & 80 GB / cu124-torch2.4 & 80 GB / cu124-torch2.4 \\
\bottomrule
\end{tabularx}
\caption{Configuration fields retained from the original profile, pending
verification. The six-character labels are short placeholders rather than full
SHA-256 digests and cannot establish reproducibility. At 81 frames and 16 fps,
nominal shot duration is 5.0625 seconds.}
\label{tab:config}
\end{table*}

\section{API Extensions and Full Resource Accounting}
\label{app:extensions}
Commercial APIs are optional transfer experiments and are excluded from the
main results. We propose a fifty-case stratified subset comparing direct
generation, global repair, and Contract2Video with one fixed provider model.
Possible providers include a Wan video endpoint or a Veo endpoint, subject to
documented availability and access at execution time. Endpoint identifiers,
versions, region, native duration, resolution, watermark policy, retries, and
pricing date are recorded. The API model is not silently substituted for an
open-weight main model, and a provider's inaccessible seed control is disclosed.

An optional audio extension uses a separately specified voice or audio model
and measures transcript adherence, speaker assignment, and audiovisual timing.
The same voice assets are supplied to every comparator. Main visual-only
renderers do not receive credit for native speech generation. ASR agreement
alone is insufficient to establish correct speaker identity or synchronization.
Audio data, consent, licensing, model costs, and human criteria are disclosed
independently of the main visual suite.

\subsection{Exact Accounting Definitions}
Let $\mathcal{J}$ contain every attempted render, including failed and retried
calls. Let $d_j^{\mathrm{bill}}$ be the provider's billable duration and $p_j$
its price per video second. Duration is rounded according to the provider's
actual billing rule, not the requested duration. Video-API cost is
\begin{equation}
C_{\mathrm{video}} = \sum_{j\in\mathcal{J}} p_jd_j^{\mathrm{bill}}
 + C_{\mathrm{fixed}} + C_{\mathrm{retry}},
\end{equation}
where the last two terms are used only for separately billed charges not
already included in the sum. They are never double-counted. Provider refunds
are recorded as explicit negative ledger items.

For model call $q$, let $t_q^{\mathrm{in}}$, $t_q^{\mathrm{out}}$, and
$t_q^{\mathrm{cache}}$ be the provider-reported billed token categories and
$p_q^{\mathrm{in}}$, $p_q^{\mathrm{out}}$, and $p_q^{\mathrm{cache}}$ their
prices per million tokens. Then
\begin{align}
C_{\mathrm{agent}}=10^{-6}\sum_q (&p_q^{\mathrm{in}}t_q^{\mathrm{in}}
 + p_q^{\mathrm{out}}t_q^{\mathrm{out}} \nonumber\\
 &+p_q^{\mathrm{cache}}t_q^{\mathrm{cache}}).
\end{align}
Input categories are mutually exclusive. Cached input is not also counted as
uncached input. Video- or image-input charges not denominated in tokens receive
separate ledger fields. Local controller calls contribute accelerator time,
not fictitious API charges.

For device $g$, measured allocation intervals $\mathcal{I}_g$ give
\begin{align}
H_{\mathrm{GPU}} &= \frac{1}{3600}
\sum_{g=1}^{8}\sum_{[a,b]\in\mathcal{I}_g}(b-a), \\
C_{\mathrm{GPU}}&=\sum_g r_gH_g.
\end{align}
where $r_g$ is the actual rental rate per GPU-hour. If the node is rented as an
indivisible eight-GPU unit, the invoice is computed from the node allocation
window and node rate instead. Unused reserved time is still charged. We report
both active-device time and billed allocation time. Owned hardware has no
pretended cloud invoice: actual energy cost and an explicitly labeled rental
equivalent may be reported separately.

The complete project cost is
\begin{align}
C_{\mathrm{total}} ={}& C_{\mathrm{GPU}} + C_{\mathrm{video}}
+C_{\mathrm{agent}} \nonumber\\
&+C_{\mathrm{audio}}+C_{\mathrm{storage}}+C_{\mathrm{egress}} \nonumber\\
&+C_{\mathrm{annotation}}+C_{\mathrm{other}}.
\end{align}
Human annotation is not omitted from the project total. Pilot, calibration,
development, failed runs, and final evaluation are distinguished from the cost
of a single deployed case. Taxes, credits, and currency conversion are stated
separately using dated invoice information. No previous conversational price
estimate is treated as a measurement or an authoritative current quote.

The optional API ledger is given in Table~\ref{tab:api-ledger}; the complete
phase-level total and per-controller efficiency are reported in
Tables~\ref{tab:total-cost} and~\ref{tab:efficiency}.

\begin{table*}[t]
\centering\tablefont
\setlength{\tabcolsep}{2pt}
\renewcommand{\arraystretch}{1.15}
\begin{tabularx}{\linewidth}{@{}L{0.15\linewidth}L{0.20\linewidth}*{6}{Y}@{}}
\toprule
\rowgreenmid
\thead{API / service} & \thead{Model / version} & \thead{Pricing date} & \thead{Billable units} & \thead{Unit price} & \thead{Gross cost} & \thead{Credits} & \thead{Net cost} \\
\midrule
\rowgreendeep
Video generation & DashScope Wan2.1-API & 2026-03-01 & 420 calls & \$0.08/call & \$33.60 & --- & \$33.60 \\
\rowgreenlight
Text controller & Local Qwen3-8B & --- & 1.84M tok & \$0.00 & \$0.00 & --- & \$0.00 \\
\rowgreenpale
Visual observer & Local Qwen2.5-VL-7B & --- & 3.12M tok & \$0.00 & \$0.00 & --- & \$0.00 \\
\rowgreenlight
Audio / speech extension & CosyVoice-v1 API & 2026-03-01 & 180 min & \$0.01/min & \$1.80 & --- & \$1.80 \\
\rowgreenpale
Storage / egress & OSS standard & 2026-03-01 & 450 GB & \$0.02/GB & \$9.00 & --- & \$9.00 \\
\bottomrule
\end{tabularx}
\caption{Unverified price and unit ledger for the optional extensions. The
main experimental comparison has no model-API dependence. Local controller and
observer calls are recorded at zero marginal API cost; provider invoices are
required to verify the listed prices.}
\label{tab:api-ledger}
\end{table*}

\begin{table*}[t]
\centering\tablefont
\setlength{\tabcolsep}{2pt}
\renewcommand{\arraystretch}{1.15}
\begin{tabularx}{\linewidth}{@{}L{0.29\linewidth}Y Y Y Y Y Y Y@{}}
\toprule
\rowgreenmid
\thead{Experimental phase} & \thead{Phase cases} & \thead{Render attempts} & \thead{GPU-hours} & \thead{Billed node-hours} & \thead{API cost} & \thead{Human cost} & \thead{Total cost} \\
\midrule
\rowgreenpale
Pilot and profiling & 20 & 168 & 14.5 & 2.0 & \$0.00 & \$0.00 & \$32.00 \\
\rowgreenlight
Development and model selection & 60 & 512 & 38.0 & 5.0 & \$0.00 & \$0.00 & \$80.00 \\
\rowgreenpale
Calibration and certification & 180 & 1,440 & 92.0 & 12.0 & \$0.00 & \$360.00 & \$552.00 \\
\rowgreenlight
Main controller comparison & 300 & 4,740 & 184.2 & 24.0 & \$0.00 & \$0.00 & \$384.00 \\
\rowgreenpale
Component ablations & 300 & 16,800 & 320.0 & 42.0 & \$0.00 & \$0.00 & \$672.00 \\
\rowgreenlight
Cross-renderer transfer & 100 & 1,620 & 54.2 & 7.0 & \$0.00 & \$0.00 & \$112.00 \\
\rowgreenpale
Offline automated evaluation & 300 & --- & 26.0 & 3.5 & \$0.00 & \$0.00 & \$56.00 \\
\rowgreenlight
Human evaluation & 150 & --- & --- & --- & \$0.00 & \$900.00 & \$900.00 \\
\rowgreenpale
Optional API / audio extensions & 50 & 420 & --- & --- & \$35.40 & \$0.00 & \$35.40 \\
\rowgreenlight
Failed runs and unplanned retries & --- & 84 & 8.5 & 1.2 & \$0.00 & \$0.00 & \$19.20 \\
\rowgreenpale
Storage, transfer, and other charges & --- & --- & --- & --- & \$9.00 & \$0.00 & \$9.00 \\
\midrule
\rowgreendeep
\textbf{Listed phase-record sum} & \textbf{1,460} & \textbf{25,784} & \textbf{737.4} & \textbf{96.7} & \textbf{\$44.40} & \textbf{\$1,260.00} & \textbf{\$2,851.60} \\
\bottomrule
\end{tabularx}
\caption{Arithmetic ledger of retained phase records, not a verified complete
project invoice. Phase cases sum to 1,460 and may reuse the same cases. A
unique-case total cannot be established from this table. Video and audio sum
to \$35.40, with \$9.00 storage counted once in its own row. The listed total
is \$2,851.60; method coverage and invoice status require manifest records.}
\label{tab:total-cost}
\end{table*}

\paragraph{Coverage requiring reconciliation.}
The listed 4,740 main-comparison attempts equal $7.9\times300\times2$, which
matches the Contract2Video row alone. Summing all seven printed call means at
the same case and seed counts gives approximately 35,100 attempts. The
component table implies approximately 37,440 attempts across its eight rows,
not the ledger's 16,800. The three video-API rows imply approximately 1,100
calls on 50 cases at one seed, rather than 420. These checks use rounded means
and do not replace measurements. Full coverage and cost remain unverified
until manifest-level accounting resolves the differences.


\begin{table*}[t]
\centering\tablefont
\setlength{\tabcolsep}{2pt}
\renewcommand{\arraystretch}{1.15}
\begin{tabularx}{\linewidth}{@{}L{0.23\linewidth}*{6}{Y}@{}}
\toprule
\rowgreenmid
\thead{Controller} & \thead{GPU-hours / case} & \thead{P50 latency} & \thead{P95 latency} & \thead{Calls / case} & \thead{Cost / case} & \thead{Cost / successful case} \\
\midrule
Direct & \metricbest{0.14}& \metricbest{162s}& \metricbest{175s}& \metricbest{4.0}& \metricbest{\$0.28}& \metricmid{\$2.33}\\
Best-of-$K$ & \metriclow{0.42}& \metriclow{486s}& \metricgood{512s}& \metriclow{12.0}& \metriclow{\$0.84}& \metriclow{\$4.35}\\
Single agent & \metricmid{0.31}& \metricmid{380s}& \metriclow{560s}& \metricmid{8.6}& \metricmid{\$0.62}& \metriclow{\$2.41}\\
Global repair & \metriclow{0.38}& \metriclow{445s}& \metriclow{680s}& \metriclow{10.4}& \metriclow{\$0.76}& \metricgood{\$1.76}\\
\papernamecolor{} & \metricgood{0.28}& \metricgood{330s}& \metricmid{520s}& \metricgood{7.9}& \metricgood{\$0.56}& \metricbest{\$1.04}\\
\bottomrule
\end{tabularx}
\caption{Retained inference profile with arithmetic cost checks, pending
verification. Printed cost per case equals GPU-hours per case times the implied
\$2 rate; successful-case costs use the rounded Independent All rate. The
case-versus-seed accounting and failed-call inclusion require a run manifest.
The successful-case cost is undefined when no case succeeds.}
\label{tab:efficiency}
\end{table*}

\subsection{Reproducibility Records}
Every intended run receives a unique case, method, seed, and configuration
identifier before execution. The artifact manifest records input hashes,
checkpoint revisions, contract versions, shot dependencies, generation seeds,
render attempts, asset hashes, verifier timestamp spans, decision thresholds,
budget counters, termination status, and measured runtime. Secrets, private
paths, account identifiers, and author-specific metadata are excluded from the
anonymous release. Failed or unresolved cases remain in the manifest.

The release will include the exact configuration used for every populated
table, paired case identifiers for statistics, external-evaluator commands,
human-rater instructions, and a machine-readable cost ledger. Large video
assets and model weights may be distributed separately subject to license,
but their hashes and retrieval instructions must be available. The present package retains a populated numerical profile without the run-level
provenance needed to validate it. Independent multi-seed replication and
reconciliation against original run records remain required.

\section{Additional Comparisons and Figure Templates}
\label{app:figure-templates}

\subsection{Published Workflow Comparisons}
We distinguish an original-system reproduction from a controller adaptation
that replaces the original renderer with the shared Wan backbone. Only the
latter isolates orchestration at a matched rendering budget. The published
systems below are targeted comparators. Their task compatibility, public
implementation revision, and deviations are documented for the supplied
fixed-profile comparison. VideoAgent addresses robotic
video planning and is discussed as related work rather than silently treated
as an equivalent movie-production baseline.

The matched-budget workflow comparison is listed in Table~\ref{tab:published-workflows}.

\begin{table*}[t]
\centering\tablefont
\setlength{\tabcolsep}{2pt}
\renewcommand{\arraystretch}{1.15}
\begin{tabularx}{\linewidth}{@{}L{0.23\linewidth}L{0.31\linewidth}*{4}{Y}@{}}
\toprule
\rowgreenmid
\thead{Workflow} & \thead{Evaluation condition} & \thead{HCS} & \thead{All} & \thead{Human win} & \thead{GPU-hours / case} \\
\midrule
VideoDirectorGPT & Shared-renderer adaptation & \metriclow{62.4\%}& \metriclow{26.8\%}& \metricmid{29.5\%}& \metricgood{0.32}\\
VideoStudio & Original-system diagnostic & \metriclow{58.1\%}& \metriclow{21.0\%}& \metriclow{24.0\%}& \metriclow{0.35}\\
VideoGen-of-Thought & Shared-renderer adaptation & \metricmid{66.8\%}& \metricmid{33.5\%}& \metricgood{36.0\%}& \metricmid{0.34}\\
MovieAgent & Shared-renderer adaptation & \metricgood{69.2\%}& \metricgood{38.0\%}& \metricbest{39.5\%}& \metriclow{0.36}\\
\papernamecolor{} & Shared renderer and budget & \metricbest{78.6\%}& \metricbest{54.0\%}& --- & \metricbest{0.28}\\
\bottomrule
\end{tabularx}
\caption{Published-workflow comparison. An adaptation is explicitly
identified and is not presented as the original authors' reported system.
Incompatible or unavailable systems are reported with reasons rather than
replaced by a weak surrogate bearing the same name.}
\label{tab:published-workflows}
\end{table*}

\subsection{Optional Provider and Audio Outcomes}
These analyses use separate manifests and never enter the open-weight primary
table. Provider requests use an identical endpoint, duration, resolution, and
billing configuration across controllers. For native audiovisual endpoints,
sound quality and semantic alignment are assessed independently of the visual
metrics. No outcome is extrapolated from an API price or a model description.

The optional provider and audio outcomes are summarized in Table~\ref{tab:api-outcomes}.

\begin{table*}[t]
\centering\tablefont
\setlength{\tabcolsep}{2pt}
\renewcommand{\arraystretch}{1.15}
\begin{tabularx}{\linewidth}{@{}L{0.16\linewidth}L{0.31\linewidth}*{5}{Y}@{}}
\toprule
\rowgreenmid
\thead{Setting} & \thead{Controller} & \thead{HCS} & \thead{All} & \thead{Calls / case} & \thead{Cost / case} & \thead{Failures} \\
\midrule
Video API & Direct & \metriclow{52.4\%}& \metriclow{16.0\%}& \metricbest{4.0}& \metricbest{\$0.32}& \metricbest{0.0\%}\\
Video API & Global repair & \metricmid{74.0\%}& \metricmid{48.0\%}& \metriclow{10.2}& \metriclow{\$0.82}& \metriclow{2.0\%}\\
Video API & \papernamecolor{} & \metricbest{81.2\%}& \metricbest{58.0\%}& \metricgood{7.8}& \metricgood{\$0.62}& \metricbest{0.0\%}\\
Audio extension & Without audio specialist & \metriclow{64.2\%}& \metriclow{31.0\%}& \metricbest{8.2}& \metricbest{\$0.66}& \metriclow{2.0\%}\\
Audio extension & \papernamecolor{} audio specialist & \metricbest{79.5\%}& \metricbest{52.0\%}& \metriclow{8.4}& \metriclow{\$0.67}& \metricbest{0.0\%}\\
\bottomrule
\end{tabularx}
\caption{Retained optional-extension summaries, pending verification. Cost/case
is the video-call subtotal from printed calls at \$0.08, rounded to two
decimals; it excludes the \$1.80 audio charge, whose per-case allocation
requires a manifest. The visual and audio constraint inventories must not be
pooled.}
\label{tab:api-outcomes}
\end{table*}

\subsection{Visual Analysis Placeholders}
The figures below are placeholders for the final curves. Any replacement uses
all intended cases and paired case-level uncertainty. Qualitative examples
follow a prespecified selection rule and include an unresolved failure. Figures
must not use synthetic data that could be mistaken for experimental observations.
The planned quality--budget, verifier-calibration, and repair examples are
shown in Figures~\ref{fig:budget},~\ref{fig:calibration}, and~\ref{fig:repair-example}.

\begin{figure*}[t]
\centering
\fbox{\parbox[c][0.21\textheight][c]{0.95\linewidth}{\centering
\textbf{Quality--budget curve placeholder}\\[8pt]
Left: independent HCS versus render-call budget\\[5pt]
Right: complete-case success versus measured accelerator time\\[5pt]
Paired case-bootstrap intervals. Shared renderer and prompt set}}
\caption{Planned quality--budget analysis. Report the initial four-shot
output and subsequent checkpoints up to the twelve-attempt ceiling. The
figure will expose whether any gain is explained by additional rendering.}
\label{fig:budget}
\end{figure*}

\begin{figure*}[t]
\centering
\fbox{\parbox[c][0.19\textheight][c]{0.95\linewidth}{\centering
\textbf{Verifier calibration placeholder}\\[8pt]
Left: selective error versus coverage\\[5pt]
Right: independently measured false acceptance by repair round\\[5pt]
Separate support and contradiction thresholds. Unknown judgments retained}}
\caption{Planned calibration diagnostic. Fixed-split certification and
trajectory-level empirical error are shown separately. The curve cannot
substitute for a formal guarantee under adaptive distribution shift.}
\label{fig:calibration}
\end{figure*}

\begin{figure*}[t]
\centering
\fbox{\parbox[c][0.25\textheight][c]{0.95\linewidth}{\centering
\textbf{Four-shot repair example placeholder}\\[8pt]
Row 1: request, hard constraints, and dependency graph\\[5pt]
Row 2: initial shots with timestamped failure evidence\\[5pt]
Row 3: selected repair closure and unchanged-asset hashes\\[5pt]
Row 4: final shots, independent verdicts, and a remaining failure}}
\caption{Planned qualitative comparison of local and global repair. The
same prompts, seeds, and render-call budget are displayed. Unchanged assets
are identified by hashes. Perceptual or semantic correctness requires
independent judgments.}
\label{fig:repair-example}
\end{figure*}


\end{document}
